\documentclass[11pt,a4paper]{ctexart}
\usepackage[margin=2.2cm]{geometry}
\usepackage{booktabs}
\usepackage{longtable}
\usepackage{array}
\usepackage{hyperref}
\hypersetup{colorlinks=true,linkcolor=blue,urlcolor=blue}
\setlength{\emergencystretch}{3em}
\sloppy

\title{Tri-EvoBench 周度工作进展报告}
\author{项目组}
\date{2026 年 8 月 28 日}

\begin{document}
\maketitle
\pagestyle{plain}

\section{本周结论}

本周完成了 Hermes L1-Life v2 的 Docker 端到端链路修复与两条 60 回合全量复测。最终使用 GLM-5.3-Flash 完成 120 个 turns；两条流均正常写出 summary，judge 请求无 HTTP 错误，平均 turn 分数分别为 0.1363 和 0.1524。为保证强制思考模型的稳定性，本次临时全量配置采用最低思考档并关闭额外后台 review；evolution barrier 仍逐回合执行并记录为 skipped。

上周已经以 LLM 跑通基础 L1 与 L2 路径。本周工作的重点不是重复 smoke test，而是消除 Hermes L1 适配、模拟副作用、跨轮评判和工作区上下文中的实际缺口，使结果可复现、可审计。

\section{上周基线状态}

\begin{longtable}{p{0.28\linewidth}p{0.24\linewidth}p{0.40\linewidth}}
\toprule
项目 & 已有结果 & 说明 \\
\midrule
基础 L1 / L2 链路 & 已跑通 & 已验证 Docker、LLM 连接、评测 Runner 与 trace 基础链路。\\
Trace2Skill & 5/5 通过，约 664 tokens & 真实 Docker + DeepSeek 回归记录位于 \url{runs/merge-trace2skill-real-fixed/}。\\
SkillOpt & 5/5 通过，约 4,644 tokens & 真实 Docker + DeepSeek 回归记录位于 \url{runs/merge-skillopt-real/}。\\
Hermes L2 & SRA E001/E002 通过，约 286,814 runtime tokens / 13 calls & 历史真实 L2 结果位于 \url{runs/merge-hermes-l2-real/}。\\
\bottomrule
\end{longtable}

\section{本周完成的工程工作}

\subsection{L1 只读工作区上下文}

L1 保持禁用 Hermes 原生工具，所有副作用仍只能经由 Tri-Evo 协议交给 Runner 执行。同时，桥接层受限地注入 \texttt{docs/} 与 \texttt{data/} 下的 UTF-8 文本快照，排除 \texttt{.git}、\texttt{state}、\texttt{iterations}、缓存与过大文件。该机制解决了分析、润色与翻译任务无法读取基准输入的问题，而不向 agent 开放原生文件写入能力。

\subsection{L1 协议与副作用修复}

\begin{itemize}
  \item L1 不再依赖 ``available\_tools 非空'' 判断模式；没有模拟工具的只读回合也会关闭原生工具并采用 L1 工作上限。
  \item 新增 \texttt{TRIEVO\_ACTION} 的 \texttt{write\_file} 解析；文件创建可将完整内容真正写入 sandbox。
  \item 对携带完整内容的既有 \texttt{file.create\_file} 协议调用进行兼容归一化，形成可审计的 \texttt{write\_file} action。
  \item 协议 prompt 明确规定：澄清回合不得发起依赖该澄清的工具调用；日历调用不得缺失 \texttt{calendar\_id}、\texttt{title}、\texttt{start}、\texttt{end}。
\end{itemize}

\subsection{跨回合状态与评判修复}

\begin{itemize}
  \item 已审计的 \texttt{state\_update} 会作为只读持久事实提供给后续 L1 回合，避免模型丢失已确认的日历 ID、输出语言或项目事实。
  \item LLM judge 现可取得当前 Turn 已发生的 interaction 轨迹，包括先前的澄清文本、事件、工具调用与状态更新；不再只看最终一句回复而误判已发生的澄清。
  \item 保留 public-trajectory 边界：judge 与 evolution 只接触公开任务、公开事件、公开工具结果和公开状态，不读取私有 oracle 内容。
\end{itemize}

\section{最终 L1-Life v2 真实回归结果}

运行配置：Docker 隔离；Hermes 镜像 \texttt{trievo-hermes-merged:0.14}；GLM-5.3-Flash（Z.AI 兼容接口）；每条流 60 turns；总 token 上限 12,000,000；\texttt{evolution.policy=disabled}；judge 使用 \texttt{reasoning\_effort=low} 与 1,024 输出上限。

\begin{center}
\renewcommand{\arraystretch}{1.15}
\begin{tabular}{lrrrrr}
\toprule
流 & turns & 交互次数 & API calls & 平均分 & 通过率 \\
\midrule
stream-1 & 60 & 62 & 131 & 0.1363 & 0\% \\
stream-2 & 60 & 62 & 139 & 0.1524 & 0\% \\
\midrule
合计 & 120 & 124 & 270 & 0.1444 & 0\% \\
\bottomrule
\end{tabular}
\end{center}

两条最终流分别消耗 2,323,736 与 2,501,417 tokens，共 270 次 API calls。所有 120 个 turns 均完成；0\% 通过率反映 GLM 在当前 L1-Life 协议下未达到严格满分阈值，不是运行中断或 judge 请求失败。此前出现的 HTTP 400/1210 已通过 GLM 强制思考参数映射与 judge 输出上限修复。

完整 summary 和 human-readable trace 位于：\url{runs/l1-life-hermes-glm-final2/stream-1/l1_life_stream_1/} 与 \url{runs/l1-life-hermes-glm-final2/stream-2/l1_life_stream_2/}。

\section{质量保证与提交}

\begin{itemize}
  \item 本地单元测试：49 项全部通过。
  \item 已执行 Python 编译检查与 \texttt{git diff --check}。
  \item 本周核心提交：\texttt{4efeaf4 Support GLM Flash judge and Hermes Docker credentials}（包含 GLM 思考参数、judge 上限与 Docker 环境透传修复）。
  \item API 密钥未写入仓库、日志、报告或提交信息；临时评测 override 与 \texttt{runs/} 结果未纳入提交。
\end{itemize}

\section{后续工作}

L1-Life v2 的两条全量链路已跑通且结果可审计。下一阶段应在相同 GLM 配置下评估 Hermes L2、SkillOpt 与 Trace2Skill，并针对 L1-Life 的低通过率改进协议提示与 agent 的结构化 JSON 产出。

\end{document}
